% Created 2026-02-19 Thu 17:35
% Intended LaTeX compiler: pdflatex
\documentclass[11pt]{article}
\usepackage[utf8]{inputenc}
\usepackage[T1]{fontenc}
\usepackage{graphicx}
\usepackage{longtable}
\usepackage{wrapfig}
\usepackage{rotating}
\usepackage[normalem]{ulem}
\usepackage{amsmath}
\usepackage{amssymb}
\usepackage{capt-of}
\usepackage{hyperref}
\date{\today}
\title{}
\hypersetup{
 pdfauthor={},
 pdftitle={},
 pdfkeywords={},
 pdfsubject={},
 pdfcreator={Emacs 30.2 (Org mode 9.7.11)}, 
 pdflang={English}}
\begin{document}

\tableofcontents

\section{PD2 writeup}
\label{sec:org465eae0}
\subsection{constants}
\label{sec:org3444244}
i made the constants such that the ALU has its own set of codes
that corresponds with the funct3 and funct7. all the opcodes have
labeled funct3s and the different types of each opcodes are made
seperate such as the different subtypes of I type.
\subsection{igen}
\label{sec:orgd562f32}
the igen is completly combonational and always gives and output
unless the opcode is R type. the immediate is muxed with the
opcode.
\subsection{decode}
\label{sec:org0a299f3}
decode is made to load on the negedge of clock. it is also
completely blocking as there is no "shift register" style logic
required by the design.
\subsection{control}
\label{sec:org398f2f5}
the control is simply a mux for the different opcodes and the
pcsel is changed from PC\textsubscript{ADD4} to PC\textsubscript{IMM} in branch and jump and the
branch opcode needs the condition from the further instruction to
be true to then change from the default PC\textsubscript{ADD4}.
\end{document}
